
\label{subsec:architecture}

The major architectural modifications of \MethodName{} model compared to previous \Pizf{} and \Pizs{} models include the use of the history vision encoder from \texttt{MEM} \cite{torne2026mem} and visual subgoal images in the context.
The model takes as input up to four camera images (front view, two wrist views, and optionally rear view), each with up to six history frames, and up to three subgoal images (omitting the rear view). The history frames are processed through the vision encoder and compressed to the same number of tokens as a single frame; subgoal images are processed through the same encoder. Both the camera observations and subgoal images are first resized to 448x448 pixels.
For sampling history frames we use a stride of 1 second, 
and the entire history frames are dropped out entirely with probability 0.3. The rear view image (when available) is dropped out with probability 0.3 as well.

We employ a block-causal masking scheme, such that the observation tokens and the subgoal image tokens use bidirectional attention within themselves, and goal-image tokens can additionally attend the observations. The following text tokens use causal attention
(see attention mask visualization in the appendix). We also feed the proprioceptive state $\bq_t$ (including the history states) of the robot into the model backbone. Unlike \Pizs{} that uses discretized text tokens to represent $\bq_t$, \MethodName{} follows \texttt{MEM} and embeds the state using a linear projection that maps the state dimension to the backbone dimension. Each history state is treated as an individual token; if the history frame is dropped out, the corresponding state token is masked out as well.

The more lightweight ``action expert'' is a 860M-parameter transformer that is trained to predict continuous actions using flow matching objective. We use adaptive RMSNorm to inject timestep information for flow matching. The number of action tokens processed by the action expert is fixed at 50, representing an action chunk of 50 steps. The 50 tokens attend bidirectionally to each other and can also attend to the VLM backbone activations.

\MethodName{} also employs the training-time version of real-time action chunking (RTC)~\cite{black2025real,black2025ttrtc} for generating smooth action trajectories in the presence of inference delay. During training, we simulate delays of 0 to 12 timesteps, corresponding to a maximum inference latency of 240ms on a 50Hz robot.

